\documentclass[10pt,a4paper]{amsart}
\usepackage[margin=19mm]{geometry}
\usepackage{amsmath,amssymb,mathtools}
\usepackage[scr=rsfs,cal=euler]{mathalfa}
\usepackage{xfrac}
\usepackage[unicode,hidelinks,bookmarks=false]{hyperref}
\newcommand{\ie}{i.e.\ }
\newcommand{\cf}{cf.\ }
\newcommand{\resp}{respectively\ }
\newcommand{\Subsection}{\subsection}
\providecommand{\nobreakdash}{\nobreak\hbox{-}}
\setlength{\emergencystretch}{2em}
\multlinegap0pt
\usepackage{amssymb}
\usepackage[normalem]{ulem}
\usepackage{xcolor}
\usepackage{enumerate}
\newtheorem{proposition}{Proposition}
\newtheorem{lemma}{Lemma}
\def\HS{\mathop{\mathrm{HS}}\nolimits}
\newcommand{\R}{\ensuremath{\mathbb{R}}}
\def\Ric{\mathop{\mathrm{Ric}}\nolimits}
\newcommand{\dd}{\mathrm{d}}
\def\div{\mathop{\mathrm{div}}\nolimits}
\newcommand{\e}{\mathrm{e}}
\newcommand{\fm}{\ensuremath{\mathfrak{m}}}
\def\trace{\mathop{\mathrm{trace}}\nolimits}
\title{Splitting theorems for weighted Finsler spacetimes\\ via the $p$-d'Alembertian: beyond the Berwald case}
\author{Erasmo CAPONIO \and Argam OHANYAN \and Shin-ichi OHTA}
\date{Source: arXiv:2412.20783v2 (CC BY 4.0)}
\begin{document}
\maketitle

\noindent\emph{Source and licence.} Splitting theorems for weighted Finsler spacetimes\\ via the $p$-d'Alembertian: beyond the Berwald case. Version arXiv:2412.20783v2 (CC BY 4.0), distributed by arXiv under the Creative Commons Attribution 4.0 International licence, http://creativecommons.org/licenses/by/4.0/, as recorded in the arXiv licence metadata for this exact version and shown on the abstract page https://arxiv.org/abs/2412.20783v2. Full licence notice: \texttt{licenses/arxiv-2412.20783-CC-BY-4.0.txt}.\par\smallskip

\noindent\textit{Editorial scope.} Counted unit: the article's lemma lm:Wylie (source0.tex lines 1213--1223). The article's complete original proof of it is delivered with it (source0.tex lines 1225--1252, one \texttt{proof} environment of 28 lines). No mathematics is written, repaired or re-derived: the statement and the proof are the article's own lines, in the article's own notation, and no retained line is substituted or re-flowed.  Retained uncounted as the source-local context the proof needs: lines 377--379; lines 1167--1171; lines 1192--1200; lines 1195--1198.  Nothing was omitted from any retained span, so the package carries no omission record.  No retained line contains a citation, so the wrapper carries no bibliography and no bibliography entry.  No retained line contains a figure environment, an image-inclusion command or drawing code, so no image is delivered and the package declares no asset dependency.  Editorial formatting only, in addition: an AMS \texttt{amsart} preamble that restates the article's own declarations and macros used by the retained lines, namely the environments \texttt{proposition}, \texttt{lemma} on separate counters, together with the standard packages those lines need.  The retained environments print in this document as Proposition 1, Lemma 1 in order of appearance, whereas the article prints the same items with the numbering of its own full document.  The complete native source of the article as distributed in the arXiv e-print for this exact version is bundled unchanged as \texttt{sources/170891/source0.tex}.
\begin{equation}\label{eq:g_X}
\frac{\dd}{\dd t}\bigl[ g_X(X,Y) \bigr] =g_X(D^X_{\dot{\eta}}X,Y) +g_X(X,D^X_{\dot{\eta}}Y)
\end{equation}

\begin{equation}\label{eq:q-HS}
\trace\bigl[ B(0)^2 \bigr]
=\sum_{i,j=1}^n g_{\nabla h}(\nabla^2_{e_i} h,e_j)^2 g_{\nabla h}(e_i,e_i) g_{\nabla h}(e_j,e_j)
=: \HS_{\nabla h}(\nabla^2 h),
\end{equation}

\begin{proposition}[Bochner-type identity]\label{pr:Boch}
Let $(M,L,\fm)$ be a weighted Finsler spacetime.
For $h \in C^3(U)$ on an open set $U \subset M$ such that $-h$ is temporal, we have
\begin{equation}\label{eq:Boch}
\div_\fm \biggl( \nabla^{\nabla h} \biggl[ \frac{F(\nabla h)^2}{2} \biggr] \biggr)
+\dd(\square_\fm h)(\nabla h) +\Ric_\infty (\nabla h) +\HS_{\nabla h}(\nabla^2 h) =0
\end{equation}
pointwise on $U$.
\end{proposition}

\begin{equation}
\div_\fm \biggl( \nabla^{\nabla h} \biggl[ \frac{F(\nabla h)^2}{2} \biggr] \biggr)
+\dd(\square_\fm h)(\nabla h) +\Ric_\infty (\nabla h) +\HS_{\nabla h}(\nabla^2 h) =0
\end{equation}

\begin{lemma}[$\Ric_0 \ge 0$ case]\label{lm:Wylie}
Let $(M,L,\fm)$ be a weighted Finsler spacetime with $\Ric_0 \ge 0$ in timelike directions.
Then, for $h \in C^3(U)$ as in Proposition~$\ref{pr:Boch}$, we have
\begin{equation}\label{eq:Wylie}
w\div_\fm \biggl( \nabla^{\nabla h} \biggl[ \frac{F(\nabla h)^2}{2} \biggr] \biggr)
+\dd(w\square_\fm h)(\nabla h) +w \frac{(\square_\fm h)^2}{n}
\le \frac{w}{2F(\nabla h)^2} \sum_{\alpha=2}^n \Bigl( \dd\bigl[ F(\nabla h)^2 \bigr](e_\alpha) \Bigr)^2
\end{equation}
at any $x \in U$, where $w(\zeta(t)):=\e^{2\psi_\fm(\dot{\zeta}(t))/n}$ for $\zeta(t):=\exp_x(t\nabla h(x))$ and $(e_i)_{i=1}^n$ is an orthonormal basis of $(T_xM,g_{\nabla h})$ such that $e_1=\nabla h(x)/F(\nabla h(x))$.
Moreover, in the case where $F(\nabla h)$ is constant, equality holds if and only if $\Ric_0(\nabla h(x))=0$ and $\nabla^2_v h=cv$ for some $c \in \R$ and all $v \in T_xM$.
\end{lemma}

\begin{proof}
We first observe
\[ \dd(w\square_\fm h)(\nabla h)
=w \cdot \dd(\square_\fm h)(\nabla h) +2w\frac{(\psi_\fm \circ \dot{\zeta})'(0)}{n} \cdot \square_\fm h.\]
It follows from \eqref{eq:q-HS}, the Cauchy--Schwarz inequality and \eqref{eq:g_X} that, at $x$,
\begin{align*}
\HS_{\nabla h}(\nabla^2 h)
&\ge \sum_{i=1}^n g_{\nabla h}(\nabla^2_{e_i} h,e_i)^2
-2\sum_{\alpha=2}^n g_{\nabla h}(\nabla^2_{e_\alpha} h,e_1)^2 \\
&\ge \frac{1}{n} \Biggl( -g_{\nabla h}(\nabla^2_{e_1} h,e_1) +\sum_{i=2}^n g_{\nabla h}(\nabla^2_{e_i} h,e_i) \Biggr)^2
-\frac{2}{F(\nabla h)^2} \sum_{\alpha=2}^n g_{\nabla h}(\nabla^2_{e_\alpha} h,\nabla h)^2 \\
&=\frac{(\square h)^2}{n} -\frac{1}{2F(\nabla h)^2} \sum_{\alpha=2}^n \Bigl( \dd\bigl[ F(\nabla h)^2 \bigr](e_\alpha) \Bigr)^2.
\end{align*}
Moreover, we have
\[ (\square h)^2 =\bigl( \square_\fm h +(\psi_\fm \circ \dot{\zeta})'(0) \bigr)^2
= (\square_\fm h)^2 +2(\psi_\fm \circ \dot{\zeta})'(0) \square_\fm h +(\psi_\fm \circ \dot{\zeta})'(0)^2. \]
Comparing these with \eqref{eq:Boch}, we obtain
\begin{align*}
&\dd(w\square_\fm h)(\nabla h) +w \frac{(\square_\fm h)^2}{n} -\frac{w}{2F(\nabla h)^2} \sum_{\alpha=2}^n \Bigl( \dd\bigl[ F(\nabla h)^2 \bigr](e_\alpha) \Bigr)^2 \\
&\le w \cdot \dd(\square_\fm h)(\nabla h) +w\HS_{\nabla h}(\nabla^2 h) -w\frac{(\psi_\fm \circ \dot{\zeta})'(0)^2}{n} \\
&= -w\div_\fm \biggl( \nabla^{\nabla h} \biggl[ \frac{F(\nabla h)^2}{2} \biggr] \biggr) -w\Ric_\infty(\nabla h) -w\frac{(\psi_\fm \circ \dot{\zeta})'(0)^2}{n} \\
&= -w\div_\fm \biggl( \nabla^{\nabla h} \biggl[ \frac{F(\nabla h)^2}{2} \biggr] \biggr) -w\Ric_0(\nabla h).
\end{align*}
Since $\Ric_0 \ge 0$ by assumption, this completes the proof of \eqref{eq:Wylie}.

When $F(\nabla h)$ is constant, the right-hand side of \eqref{eq:Wylie} vanishes.
Then, by the above proof, equality holds if and only if $\Ric_0(\nabla h(x))=0$ and $g_{\nabla h}(\nabla^2_{e_i} h,e_j)=c \cdot g_{\nabla h}(e_i,e_j)$ for some $c \in \R$, which means $\nabla^2_v h=cv$.
\end{proof}

\end{document}
